\njustchapter{绪论}

随着经济发展的进步，能源消耗的增加带来了严重的污染，严重威胁着人类的健康和环境安全。人类对能源的依赖与化石燃料燃烧之间的矛盾，推动了各种能源的开发利用，如太阳能、风能、潮汐能、核能等。因此，围绕能源转换、材料制备和智能制造的交叉研究逐渐成为本科毕业设计中的重要选题。相关研究不仅关注材料本身的结构和性能，也关注实验平台的可重复性、数据采集的稳定性以及文献资料的规范引用\cite{liu2021tool,zhai2021review}。

\section{课题背景及研究意义}

为了构建一种理想的电荷存储几何形状，电极材料通常需要同时具备较大的比表面积、稳定的孔道结构以及良好的导电网络。阳极氧化方法具有工艺简单、参数可控和形貌连续等特点，适合用于建立材料结构与电化学性能之间的对应关系。已有研究表明，微观结构差异会显著影响材料的电学性质\cite{chen2021microstructure}，这为通过工艺调控获得目标性能提供了依据。

本模板中的正文示例仅用于展示章节、公式、图表、交叉引用和 BibTeX 文献引用的写法。正式论文写作时，可保留这些语义命令并替换为自己的研究内容。

\section{超级电容器}

\subsection{超级电容器的储能机理}

在电化学储能装置中，电能主要通过双电层储能和赝电容储能两类方式保存。双电层电容器（Electrical Double-Layer Capacitors，EDLCs）通过非法拉第过程以静电方式储存电荷；赝电容器则通过快速、可逆的表面氧化还原反应储存电荷。对理想电容器而言，储能量可由\njusteqref{eq:energy-storage} 表示，其中 $E$ 为储能量，$C$ 为电容，$U$ 为工作电压。

\begin{equation}
  E = \frac{1}{2} C U^2
  \label{eq:energy-storage}
\end{equation}

\subsection{性能评价指标}

常用评价指标包括面积比电容、质量比电容、循环稳定性和倍率性能等。为了保证结果可比较，实验记录中需要同时说明电解液浓度、扫描速率、活性材料质量和电极有效面积。图、表和公式的交叉引用示例见第~\ref{chap:experiment-sample} 章。
